\section{A Lorentz-distance family}\label{sec:family}

Throughout this section, $\Sigma^m\subset S^{m+1}$ is smooth, closed,
connected, embedded, minimal, and not totally geodesic, with $m\ge2$.
Write $|\Sigma|=\Vol_m(\Sigma)$ and $d\mu$ for its area measure.
Choose a global unit normal $N$, with shape operator $A=-dN$.
The construction below uses the signed Euclidean distance to the cone over
$\Sigma$.  Acting on its graph by Lorentz transformations produces a family
whose area can be estimated on the fixed source $\Sigma$, including at cut
points and at infinite boosts.
The use of distance levels, compactified parameters, and a local index
deformation has antecedents in the canonical-family construction of
Marques and Neves \cite[Sections~2.3 and~10]{MarquesNeves2014}.
Here the transformation acts on the graph of Euclidean cone distance;
the complementary fillings will supply projective detection.

\begin{proposition}[The Lorentz-distance family]\label{prop:lorentz-family}
Let $p=m+3$ and
\[
 Q=\overline B^{m+2}\times[-1,1],\qquad u=(v,b).
\]
Identify opposite points of $\partial Q$, and denote the quotient map by
$q:Q\longrightarrow\mathbb {RP}^{p}$.  There is a flat-continuous map
\[
 \Phi:\mathbb {RP}^{p}\longrightarrow
 \mathcal Z_m^0(S^{m+1};\mathbb Z_2)
\]
with the following properties.
\begin{enumerate}[label=\textup{(\roman*)}]
\item The map $\Phi$ is detecting in the sense of
Definition~\ref{def:detecting}, has no concentration of mass, and satisfies
\[
 \Phi(q(0,0))=\partial[\Omega],\qquad
 \mathbf M(\Phi(q(0,0)))=|\Sigma|,
\]
where $\Omega$ is the side of $\Sigma$ opposite to $N$.
For $r=\sqrt{1-b^2}$ its masses obey
\begin{equation}\label{eq:family-mass-bounds}
 \mathbf M(\Phi(q(v,b)))\le
 \begin{cases}
  r^m|\Sigma|,& |v|<1,\\
  r^m\sigma_m,& |v|=1.
 \end{cases}
\end{equation}
In particular $|\Sigma|\ge\gamma_m>\sigma_m$, where $\gamma_m$ is the
constant in Lemma~\ref{lem:detection-gap}.
\item The central parameter is the unique maximum, in the uniform sense
that for every neighborhood $U$ of $(0,0)$ in $Q$,
\begin{equation}\label{eq:family-away-gap}
 \sup_{u\in Q\setminus U}\mathbf M(\Phi(q(u)))<|\Sigma|.
\end{equation}
\item Near $(0,0)$ these cycles are precisely the multiplicity-one
boundaries of a smooth family of domains, with embedded boundaries
parametrized by $\Sigma$.  The normal derivative at the center, in the
parameter direction $(d,t)\in\mathbb R^{m+2}\times\mathbb R$, is
\begin{equation}\label{eq:family-normal-variations}
 w_{d,t}(x)=t+\langle d,x\rangle.
\end{equation}
These variations form a $p$-dimensional subspace on which the area second
variation is negative definite.
\item If $\ind(\Sigma)>p$, there is a flat homotopy of $\Phi$, supported
in the image of a neighborhood of $(0,0)$, to a detecting map
$\widetilde\Phi$ with no concentration of mass and
\begin{equation}\label{eq:family-strict-deformation}
 \sup_{\xi\in\mathbb {RP}^{p}}\mathbf M(\widetilde\Phi(\xi))<|\Sigma|.
\end{equation}
\end{enumerate}
\end{proposition}

We first parametrize closest-point contacts and estimate every level.
The same source estimate then gives no concentration, including at
infinite boosts, and the complementary boundary fillings establish
detection. Finally, we isolate the central maximum and deform it using
an additional negative direction. The only global area input needed
to begin the construction is $|\Sigma|\ge\sigma_m$, furnished by cone
monotonicity in Proposition~\ref{prop:immersion-reduction}.

\subsection{Signed distance and closest-point contact}

Let $\mathcal C=C(\Sigma)\subset\mathbb R^{m+2}$ include its vertex.
The two sides of $\Sigma$ determine two sides of $\mathcal C$.
Let $\varphi$ be signed Euclidean distance to $\mathcal C$, positive on
the side pointed to by $N$ and zero on $\mathcal C$.
This function is homogeneous of degree one under positive dilations and
is $1$-Lipschitz.  For points on the same side, the latter assertion is
the ordinary distance inequality.  For points on opposite sides, their
joining segment meets $\mathcal C$, and the sum of their distances to
$\mathcal C$ is at most the length of that segment.

\begin{lemma}[Cone contact]\label{lem:family-cone-contact}
There is $K<1$ such that
\begin{equation}\label{eq:family-spacelike-bound}
 |\varphi(z)|\le K|z|\qquad(z\in\mathbb R^{m+2}).
\end{equation}
For $z\ne0$, put $h=\varphi(z)$.  Every closest point of $\mathcal C$ to
$z$ has the form $r_zx$, with $x\in\Sigma$ and
\begin{equation}\label{eq:family-contact}
 r_z=\sqrt{|z|^2-h^2}>0,\qquad
 z=r_zx+hN(x),\qquad r_z I-hA(x)\ge0.
\end{equation}
The final inequality is an inequality of self-adjoint endomorphisms of
$T_x\Sigma$.
\end{lemma}

\begin{proof}
Define $H(z)=\max_{x\in\Sigma}\langle z,x\rangle$.
Lemma~\ref{lem:fullness} gives $\int_\Sigma x\,d\mu=0$ and fullness.
Thus $H(z)>0$ for $z\ne0$: a nonpositive coordinate function of mean
zero would vanish identically.  Compactness gives
$c=\min_{|z|=1}H(z)>0$.
Minimizing first along each ray of the cone gives
\[
 \dist(z,\mathcal C)^2=|z|^2-H(z)^2
 \le(1-c^2)|z|^2.
\]
This proves \eqref{eq:family-spacelike-bound} with
$K=\sqrt{1-c^2}$ and excludes the vertex as a closest point to nonzero
$z$.

A closest ray is determined by a maximum $x$ of
$\langle z,\cdot\rangle$, and its radial coordinate is $H(z)$.
At this maximum the tangential component of $z$ vanishes, so
$z=H(z)x+\widetilde hN(x)$, where $|\widetilde h|=|h|$.
The segment from $z$ to its closest point meets no earlier cone point.
Its side near contact therefore agrees with the side of $z$, proving
$\widetilde h=h$ with our signed-distance convention.
Finally, at this height maximum the spherical immersion formula gives
\[
 \nabla^2_\Sigma\langle z,x\rangle=-H(z)I+hA\le0.
\]
These arguments apply to each maximum separately and therefore also
apply when there are several closest points.
\end{proof}

\subsection{The transformed graph and its level parametrization}

For $v\in B^{m+2}$ write
\[
 v=\tanh\tau\,a,\qquad
 C=\cosh\tau,\quad S=\sinh\tau,\quad
 B=I+(C-1)aa^t,
\]
where $|a|=1$ and $\tau\ge0$.  At $v=0$, the expressions below are
understood by continuity; they do not depend on the choice of $a$.
Indeed, $Sa=Cv$ and
$(C-1)aa^t=C^2vv^t/(C+1)$ are smooth in $v$.
Use the Lorentz transformation
\begin{equation}\label{eq:family-boost}
 z=By+bSa,\qquad h=Cb+S\langle a,y\rangle,
 \qquad |z|^2-h^2=|y|^2-b^2.
\end{equation}

\begin{lemma}[Global graph and compactification]\label{lem:family-graph}
For each $v\in B^{m+2}$ the inverse image under
\eqref{eq:family-boost} of $h=\varphi(z)$ is the graph of a continuous
function $b=\psi_v(y)$ on all of $\mathbb R^{m+2}$.
The functions depend continuously on $(v,y)$.
On $S^{m+1}$ they satisfy $|\psi_v|<1$ and extend jointly continuously
to $\overline B^{m+2}\times S^{m+1}$ by
\begin{equation}\label{eq:family-boundary-graph}
 \psi_a(y)=-\langle a,y\rangle\qquad(|a|=1).
\end{equation}
\end{lemma}

\begin{proof}
For fixed $v,y$, consider
\[
 F_{v,y}(b)=Cb+S\langle a,y\rangle-\varphi(By+bSa).
\]
The Lipschitz property of $\varphi$ gives, when $b_2>b_1$,
\[
 F_{v,y}(b_2)-F_{v,y}(b_1)\ge(C-S)(b_2-b_1)>0.
\]
The same estimate implies limits $-\infty$ and $+\infty$ at the two
ends of the real line.  Hence there is exactly one zero, and the local
positive lower bound for $C-S$ proves continuous dependence of that
zero on finite parameters and $y$.

At a graph point with $|y|=1$, \eqref{eq:family-spacelike-bound} and
Lorentz invariance give
\begin{equation}\label{eq:family-graph-estimate}
 1-b^2=|z|^2-h^2\ge(1-K^2)|z|^2,
 \qquad
 |h|\le\frac{K}{\sqrt{1-K^2}}\sqrt{1-b^2}.
\end{equation}
Equality $|b|=1$ would force $z=h=0$, contrary to $|y|=1$ and
invertibility of the boost.  Thus $|b|<1$.
Since $b+\langle v,y\rangle=h/C$, we also obtain
\begin{equation}\label{eq:family-uniform-compactification}
 \sup_{|y|=1}|\psi_v(y)+\langle v,y\rangle|
 \le\frac{K}{C\sqrt{1-K^2}}.
\end{equation}
This tends to zero as $|v|\to1$, uniformly in $y$ and in the boost
direction, proving the claimed extension.
\end{proof}

Fix now $|v|<1$ and $-1<b<1$, and put $r=\sqrt{1-b^2}$.
For $x\in\Sigma$ define
\begin{equation}\label{eq:family-level-formulas}
 \begin{gathered}
 f=\langle a,x\rangle,\qquad g=\langle a,N\rangle,
 \qquad D=C-Sg>0,\\
 N'=\frac{BN-Sa}{D},\qquad
 h_{v,b}(x)=\frac{b+Srf}{D},\\
 Y_{v,b}(x)=r(Bx+SfN')+bN'.
 \end{gathered}
\end{equation}
These are obtained by applying the inverse boost
\[
 y=Bz-Sh a,\qquad b=Ch-S\langle a,z\rangle
\]
to $z=rx+hN$.  In fact the second equation becomes $b=Dh-Srf$,
which gives precisely the displayed $h_{v,b}$ and $Y_{v,b}$.

The differential of this map has an especially useful form.
For $w\perp N(x)$ set
\[
 Ow=Bw+S\langle a,w\rangle N'.
\]
The identities $B^2=I+S^2aa^t$ and
\[
 |BN-Sa|^2=D^2,\qquad
 \langle Bw,N'\rangle=-S\langle a,w\rangle
 \quad(w\perp N)
\]
show that $|N'|=1$ and that
$O:N^\perp\longrightarrow(N')^\perp$ is a linear isometry.
In particular,
\begin{equation}\label{eq:family-orthogonal-decomposition}
 Y_{v,b}=rOx+bN'\in S^{m+1},\qquad
 \langle N',Y_{v,b}\rangle=b.
\end{equation}
The two pointwise normal planes in Figure~\ref{fig:family-normal-planes}
explain the roles of the scalars.  In
$z=rx+hN$, the ratio $h/r$ is the tangent of the signed angle from
$x$ to $z/|z|$; in \eqref{eq:family-orthogonal-decomposition}, $b$ is
the sine of the signed angle from $Ox$ to $Y_{v,b}$.
Although $r$ is fixed on a level, $h_{v,b}$ generally varies with $x$.

\begin{figure}[ht]
\centering
\begin{tikzpicture}[>=Latex,font=\small]
  \begin{scope}
    \node at (1.2,2.9) {Source contact plane};
    \draw[->] (0,0) -- (2.7,0) node[right] {$x$};
    \draw[->] (0,0) -- (0,2.5) node[left] {$N(x)$};
    \draw[dashed] (1.6,0) -- (1.6,1.8);
    \draw[thick,->,linkblue] (0,0) -- (1.6,1.8);
    \fill (1.6,1.8) circle (1.3pt);
    \node[above right] at (1.6,1.8) {$z=rx+hN$};
    \node[below] at (0.8,0) {$r$};
    \node[right] at (1.6,0.9) {$h$};
  \end{scope}
  \begin{scope}[xshift=7cm]
    \node at (1.2,2.9) {Transformed normal plane};
    \draw[->] (0,0) -- (2.7,0) node[right] {$Ox$};
    \draw[->] (0,0) -- (0,2.5) node[left] {$N'(x)$};
    \draw[gray] (2,0) arc[start angle=0,end angle=100,radius=2];
    \draw[dashed] (1.6,0) -- (1.6,1.2);
    \draw[thick,->,linkblue] (0,0) -- (1.6,1.2);
    \fill (1.6,1.2) circle (1.3pt);
    \node[above right] at (1.6,1.2) {$Y=rOx+bN'$};
    \node[below] at (0.8,0) {$r$};
    \node[right] at (1.6,0.6) {$b$};
    \node at (2.6,2.1) {$|Y|=1$};
  \end{scope}
\end{tikzpicture}
\caption{The two decompositions at a fixed closest-point contact, shown
schematically with positive normal coefficients. Lorentz invariance
makes the coefficient $r=\sqrt{1-b^2}$ the same in both planes.
The source vector $z$ need not be unit length. On a fixed level, $b$
is constant while $h=h_{v,b}(x)$ can vary. The isometry $O$ and normal
$N'$ depend on the source point; the picture does not represent a
single ambient rotation.}
\label{fig:family-normal-planes}
\end{figure}

Differentiating $N'$ and then $Y_{v,b}$, with $v,b$ fixed, gives
\[
 dN'=D^{-1}\bigl(B\,dN+S\,dg\,N'\bigr)=D^{-1}O\,dN,
\]
and consequently
\begin{equation}\label{eq:family-jacobian}
 \begin{split}
 dY_{v,b}
 &=rO\,dx+(b+rSf)D^{-1}O\,dN
   =O(rI-h_{v,b}A),\\
 J_mY_{v,b}&=\bigl|\det(rI-h_{v,b}A)\bigr|.
 \end{split}
\end{equation}
Here $J_m$ is the metric $m$-dimensional Jacobian; the isometry $O$
accounts for the entire ambient change in the differential.

\subsection{Every level has the required perimeter bound}

Define
\[
 E_{v,b}=\{y\in S^{m+1}:\psi_v(y)<b\},\qquad
 \Lambda_{v,b}=\{y\in S^{m+1}:\psi_v(y)=b\}.
\]
For finite boosts and interior levels, also define the closed source set
\[
 T_{v,b}=\{x\in\Sigma:rI-h_{v,b}(x)A(x)\ge0\}.
\]

\begin{lemma}[The restricted area bound]\label{lem:family-all-levels}
Every $E_{v,b}$, for $(v,b)\in Q$, has finite perimeter.
For $|v|<1$, $|b|<1$ and every open set $U\subset S^{m+1}$,
\begin{equation}\label{eq:family-restricted-perimeter}
 \begin{split}
 \operatorname{Per}(E_{v,b};U)
 &\le\mathcal H^m(\Lambda_{v,b}\cap U)\\
 &\le\int_{T_{v,b}\cap Y_{v,b}^{-1}(U)}
       \det(rI-h_{v,b}A)\,d\mu,
 \end{split}
\end{equation}
and on $T_{v,b}$,
\begin{equation}\label{eq:family-restricted-density}
 0\le\det(rI-h_{v,b}A)\le r^m\le1.
\end{equation}
Thus these estimates hold at every level, without an exceptional set
of values of $b$.
\end{lemma}

\begin{proof}
Take $y\in\Lambda_{v,b}$ and transform $(y,b)$ to $(z,h)$ by
\eqref{eq:family-boost}.  Lemma~\ref{lem:family-cone-contact} supplies
a closest point $r_zx$.  Lorentz invariance gives
$r_z^2=|z|^2-h^2=1-b^2$, so $r_z=r$.
The inverse formulas then show that $h=h_{v,b}(x)$,
$y=Y_{v,b}(x)$, and $x\in T_{v,b}$.  Therefore
\begin{equation}\label{eq:family-level-cover}
 \Lambda_{v,b}\subset Y_{v,b}(T_{v,b}).
\end{equation}
The nonnegative eigenvalues of $rI-h_{v,b}A$ have sum $mr$, because
$\tr A=0$.  Their arithmetic--geometric mean inequality proves
\eqref{eq:family-restricted-density}. The area formula
\cite[\S3.2.3]{Federer1969}\ \cite[\S3, pp.~299--300]{Federer1978},
with multiplicities counted in its source
integral, now yields the second inequality in
\eqref{eq:family-restricted-perimeter}.
The condition defining $T_{v,b}$ is only necessary for a closest-point
contact.  Additional points in this source set, and multiple preimages
of a level point, increase the upper bound and cause no difficulty.

In particular $\mathcal H^m(\Lambda_{v,b})<\infty$.
Since $E_{v,b}$ is open, its essential boundary is contained in its
topological boundary, which is contained in $\Lambda_{v,b}$.
The finite-essential-boundary criterion
\cite[\S4, p.~302]{Federer1978}, applied in coordinate charts, gives
finite perimeter. The intrinsic reduced-boundary representation
\cite[\S2.2]{BDFSphere} then gives
\[
 \operatorname{Per}(E_{v,b};U)
 =\mathcal H^m(\partial^*E_{v,b}\cap U)
 \le\mathcal H^m(\Lambda_{v,b}\cap U).
\]
This proves the first inequality with constant one.

For $|v|<1$ the endpoint sublevels are empty and full, respectively.
For $|v|=1$, formula~\eqref{eq:family-boundary-graph} makes them round
caps, whose boundary masses are $\sigma_m(1-b^2)^{m/2}$, with zero
mass at $b=\pm1$.  These cases also have finite perimeter.
\end{proof}

All of the levels $\Lambda_{v,b}$ have zero $(m+1)$-dimensional volume:
this follows from their finite $m$-dimensional measure in the interior;
on the boost boundary they are round spheres or singletons; and at
finite-boost endpoints they are empty.
If $(v_j,b_j)\to(v,b)$, Lemma~\ref{lem:family-graph} gives uniform
convergence of $\psi_{v_j}-b_j$ to $\psi_v-b$.  Indicators of the
sublevels consequently converge pointwise off $\Lambda_{v,b}$.
Dominated convergence proves
\begin{equation}\label{eq:family-filling-continuity}
 \Vol_{m+1}(E_{v_j,b_j}\mathbin\triangle E_{v,b})\longrightarrow0.
\end{equation}
Thus
\[
 \widehat\Phi(v,b)=\partial[E_{v,b}]
\]
is flat continuous on $Q$, since the flat distance between two of these
cycles is at most the volume of the symmetric difference of their
fillings.  Its mass equals the perimeter of the filling, so
\eqref{eq:family-mass-bounds} follows from
Lemma~\ref{lem:family-all-levels}.

\subsection{Infinite boosts and no concentration of mass}

We have obtained a flat-continuous family with the required total mass
bound. To apply the min--max theorem, we must also control its mass in
small balls uniformly over parameters, including boosts tending to
infinity. The restricted Jacobian estimate lets us do this on the one
fixed compact source \(\Sigma\).

\begin{lemma}[No concentration]\label{lem:family-no-concentration}
The family $\widehat\Phi$ satisfies
\begin{equation}\label{eq:family-no-concentration}
 \lim_{\rho\downarrow0}\;
 \sup_{(v,b)\in Q}\sup_{y\in S^{m+1}}
 \|\widehat\Phi(v,b)\|(B_\rho(y))=0.
\end{equation}
\end{lemma}

\begin{proof}
We prove the sequential form of the assertion.  Consider
$u_j=(v_j,b_j)\to(v_*,b_*)$, $y_j\to y_*$, and $\rho_j\downarrow0$.
The case $|b_*|=1$ follows directly from
\eqref{eq:family-mass-bounds}.  Parameters on $|v|=1$ are harmless as
well: an $m$-sphere of Euclidean radius $r$ in any affine
$(m+1)$-plane has at most $c_m\rho^m$ area in a ball of radius $\rho$,
uniformly in $r$ and the plane.  For $r\le2\rho$ use its total area;
for $r>2\rho$ the intersection lies in a spherical cap with angular
radius at most $\arcsin(\rho/r)$, and the same bound follows by the
cap area formula.  Spherical balls in the ambient unit sphere lie in
Euclidean balls of the same radius.  We may therefore pass to a
subsequence with $|v_j|<1$ and $|b_*|<1$.

Put $Y_j=Y_{v_j,b_j}$ and $T_j=T_{v_j,b_j}$.
The local perimeter bound reduces the desired assertion to
\begin{equation}\label{eq:family-concentration-integrands}
 \int_\Sigma I_j\,d\mu\longrightarrow0,
 \qquad
 I_j=\mathbf1_{T_j}\,
      \mathbf1_{Y_j^{-1}(B_{\rho_j}(y_j))}\,J_mY_j.
\end{equation}
The restricted bound gives $0\le I_j\le1$ on the fixed compact source.

If $|v_*|<1$, the maps $Y_j$ converge smoothly on $\Sigma$ to
$Y_*=Y_{v_*,b_*}$.  Away from the fiber $Y_*^{-1}(y_*)$, the ball
indicator in \eqref{eq:family-concentration-integrands} eventually
vanishes.  The area formula applied to the fiber gives
\[
 \int_{Y_*^{-1}(y_*)}J_mY_*\,d\mu=0,
\]
because a singleton has zero $\mathcal H^m$ measure.
Thus $J_mY_*=0$ almost everywhere on that fiber, and smooth
convergence gives $I_j\to0$ there almost everywhere as well.
Dominated convergence completes this case.

It remains to consider $v_*=a$ with $|a|=1$.
Write $a_j=v_j/|v_j|$, $t_j=S_j/C_j=|v_j|$, and
$f_j=\langle a_j,x\rangle$, $g_j=\langle a_j,N\rangle$.
The cancellation identity
\begin{equation}\label{eq:family-cancelled-map}
 Bx+SfN'=x+f\,\frac{SN-(C-1)a}{D}
\end{equation}
gives
\[
 Y_j=r_j\left[x+
 f_j\frac{t_jN-(1-C_j^{-1})a_j}{1-t_jg_j}\right]+b_jN'_j.
\]
On compact subsets of $\Sigma\setminus Z_a$, where
$Z_a=\{x:N(x)=a\}$, the denominators stay uniformly positive.
Also
\[
 N'_j=
 \frac{C_j^{-1}N+(1-C_j^{-1})g_ja_j-t_ja_j}{1-t_jg_j}
 \longrightarrow-a.
\]
The maps and all their source derivatives therefore converge locally
to
\begin{equation}\label{eq:family-infinite-map}
 Y_\infty(x)=r_*\left[x+
       \frac{\langle a,x\rangle}{1-\langle a,N\rangle}(N-a)
                 \right]-b_*a,
 \qquad r_*=\sqrt{1-b_*^2}.
\end{equation}
This computation allows the directions $a_j$ to vary.

The exceptional normal fiber $Z_a$ is source-null.
Indeed $Z_a\subset\{x:\langle a,x\rangle=0\}$, whereas
$\langle a,x\rangle$ is a nontrivial analytic function by
Lemma~\ref{lem:fullness}.  Its zero set has zero $m$-dimensional
measure.  No rank assumption on the Gauss map is involved.

The point fiber $F=Y_\infty^{-1}(y_*)$, defined off $Z_a$, may instead
have positive source measure.  On a countable cover by relatively
compact coordinate patches in $\Sigma\setminus Z_a$, the area formula
gives
\[
 \int_{F\cap K}J_mY_\infty\,d\mu=0
\]
for each such patch $K$.  Hence $J_mY_\infty=0$ almost everywhere on
$F$.  Off $F\cup Z_a$ the shrinking-ball indicator is eventually
zero; on $F$, local smooth convergence gives $J_mY_j\to0$ almost
everywhere.  We conclude that $I_j\to0$ almost everywhere on
$\Sigma$, and $0\le I_j\le1$ again permits dominated convergence.
In particular, convergence of the indicators $\mathbf1_{T_j}$ is
unnecessary.  Only the restricted Jacobians are bounded; no bound for
the unrestricted $J_mY_j$ near $Z_a$ has been used.

Finally, if \eqref{eq:family-no-concentration} failed, compactness of
$Q$ and $S^{m+1}$ would supply sequences of the above form whose ball
masses stay bounded below by a positive constant.  The sequential
result excludes them.
\end{proof}

\subsection{The projective parameter space}

On $|v|=1$, formula~\eqref{eq:family-boundary-graph} gives
\[
 E_{v,b}=\{\langle v,y\rangle>-b\},\qquad
 E_{-v,-b}=\{\langle v,y\rangle<-b\}.
\]
These fillings are complementary modulo their null common level.
On the top and bottom faces they are full and empty, respectively.
Together with the volume continuity in
\eqref{eq:family-filling-continuity}, these facts verify every
hypothesis of Lemma~\ref{lem:minmax-projective}. Therefore
$\widehat\Phi(u)=\partial[E_u]$ descends to a flat-continuous detecting
map $\Phi$ on $\mathbb {RP}^{p}$. Its image is unchanged by passage to
the quotient, so Lemma~\ref{lem:family-no-concentration} also gives no
concentration for $\Phi$.

At $(0,0)$, $\psi_0=\varphi$ on the sphere, and the level is
$\Sigma$, with sublevel $\Omega$.  Combining the upper mass bound,
Lemma~\ref{lem:family-no-concentration}, and
Lemma~\ref{lem:detection-gap} proves part~\textup{(i)} of
Proposition~\ref{prop:lorentz-family}, in particular
\begin{equation}\label{eq:family-gap}
 |\Sigma|\ge\gamma_m>\sigma_m.
\end{equation}

\subsection{Uniform area loss away from the center}

The family now detects the required cohomology class and has no
concentration. We next prove that all parameters away from the original
hypersurface lose a uniform amount of area. This isolates the maximum
that the extra negative direction will remove.

\begin{lemma}[A fixed source patch gives strict loss]
\label{lem:family-strict-loss}
For every neighborhood $U$ of $(0,0)$,
\eqref{eq:family-away-gap} holds.
\end{lemma}

\begin{proof}
Suppose otherwise.  There are $u_j=(v_j,b_j)\in Q\setminus U$ with
$\mathbf M(\widehat\Phi(u_j))\to|\Sigma|$.
The mass bound forces $b_j\to0$.
Parameters on the boost boundary have mass at most $\sigma_m$, which
is strictly smaller by \eqref{eq:family-gap}; discard those terms.
After taking a subsequence, $v_j\to v_*\ne0$.
Set $t_*=|v_*|\in(0,1]$ and $a=v_*/|v_*|$.
Discard finitely many terms so that $v_j\ne0$.
For each $j$, set $t_j=|v_j|$, $a_j=v_j/|v_j|$,
$C_j=(1-t_j^2)^{-1/2}$, $r_j=\sqrt{1-b_j^2}$,
$f_j=\langle a_j,x\rangle$, and $g_j=\langle a_j,N\rangle$.

The open set $\{A\ne0\}$ is nonempty.  The analytic function
$f=\langle a,x\rangle$ cannot vanish on that entire open set, by
fullness and connectedness.  Choose a compact source patch $K_0$ of
positive measure on which $|A|$ and $|f|$ are bounded below by
positive constants.  It is also disjoint from $\{N=a\}$, because
$N=a$ implies $f=0$.
On this fixed patch, the formula for $h_j=h_{v_j,b_j}$ becomes
\[
 h_j=\frac{b_j/C_j+t_jr_jf_j}{1-t_jg_j}
 \longrightarrow
 h_*:=\frac{t_*f}{1-t_*\langle a,N\rangle},
\]
uniformly.  Its limit is finite and nonzero on $K_0$, including when
$t_*=1$.

For a real symmetric matrix $M$, define its truncated determinant by
\[
 \mathcal D(M)=
 \begin{cases}
  \det M,&M\ge0,\\
  0,&M\not\ge0.
 \end{cases}
\]
This is continuous: every boundary point of the positive-semidefinite
cone is singular, so its determinant vanishes.
At every $x\in K_0$, the matrix $I-h_*(x)A(x)$ has trace $m$ and is
not the identity.  If it is positive semidefinite, strictness in the
arithmetic--geometric mean inequality gives determinant less than
one; if it is not, $\mathcal D$ is zero.
Compactness of $K_0$ and continuity therefore give $\delta>0$ with
\[
 \mathcal D(I-h_*A)\le1-2\delta\quad\hbox{on }K_0.
\]
For all large $j$ this yields
$\mathcal D(r_jI-h_jA)\le r_j^m-\delta$ on $K_0$.
On all of $\Sigma$, the restricted determinant is at most $r_j^m$.
Thus the area formula estimate gives
\[
 \mathbf M(\widehat\Phi(u_j))
 \le\int_\Sigma\mathcal D(r_jI-h_jA)\,d\mu
 \le r_j^m|\Sigma|-\delta\mu(K_0),
\]
contradicting convergence to $|\Sigma|$.
The proof supplies a uniform loss on a fixed portion of the source;
it makes no upper-semicontinuity assumption about area in the flat
topology.
\end{proof}

\subsection{The entire central level is a smooth boundary}

\begin{lemma}[The central tube]\label{lem:family-central-tube}
For all sufficiently small $(v,b)$, the entire level
$\Lambda_{v,b}$ is $Y_{v,b}(\Sigma)$, with $Y_{v,b}$ a smooth
embedding depending smoothly on the parameters.  Every level point
has a unique closest-point contact in the cone representation.
Moreover
\begin{equation}\label{eq:family-regular-level}
 \nabla_{S^{m+1}}\psi_v=N'-bY_{v,b},\qquad
 |\nabla_{S^{m+1}}\psi_v|=r>0
 \quad\hbox{on }\Lambda_{v,b}.
\end{equation}
Thus $\Lambda_{v,b}=\partial E_{v,b}=\partial^*E_{v,b}$ and its
boundary cycle has multiplicity one.
\end{lemma}

\begin{proof}
Use a tubular neighborhood of the smooth cone on a compact annulus
containing $\{sx:x\in\Sigma,\ 1/2\le s\le3/2\}$, constructed inside
a slightly larger annulus.  Compact embeddedness supplies a uniform
$\epsilon>0$ so that, for these central radii and $|h|<\epsilon$,
the point $sx+hN(x)$ has unique closest cone point $sx$, and signed
distance $h$.  This is a tube for the whole cone: after shrinking
$\epsilon$, points at radii outside the larger annulus cannot be
closest, and uniqueness within that annulus follows from its tubular
neighborhood.

For any actual level point $y\in\Lambda_{v,b}$, the boost gives
\[
 |h|=|Cb+S\langle a,y\rangle|\le C|b|+S\longrightarrow0
 \quad\hbox{as }(v,b)\to(0,0).
\]
Every closest cone point has radius $r=\sqrt{1-b^2}\to1$.
Thus all actual contacts lie in this one fixed tube, so each is
unique.  Conversely, for every source point $x$, the inverse formula
gives
\[
 |h_{v,b}(x)|\le\frac{|b|+Sr}{C-S}\longrightarrow0
\]
uniformly on $\Sigma$.
Consequently $z=rx+h_{v,b}(x)N(x)$ lies in the same tube and really
has signed distance $h_{v,b}(x)$.  Its transformed point
$Y_{v,b}(x)$ is therefore an actual level point.
This proves equality with the entire level and uniqueness of its
source representation, not merely existence of one nearby component.

For these parameters the matrix $rI-h_{v,b}A$ is uniformly positive
definite, so \eqref{eq:family-jacobian} gives an immersion.
Uniqueness of the contact and invertibility of the boost give
injectivity; compactness then gives an embedding.
All the formulas are smooth in $(v,b,x)$ near the center.

In the tube $d\varphi=N\cdot dz$.  Differentiating the graph relation
at one of the level points yields
\[
 0=dh-N\cdot dz
  =D\bigl(db-N'\cdot dy\bigr).
\]
The derivative with respect to $b$ is $D>0$, so the graph is smooth
near that point and its Euclidean gradient is $N'$.
Projecting onto the tangent space of the unit sphere and using
\eqref{eq:family-orthogonal-decomposition} proves
\eqref{eq:family-regular-level}.
Each level point therefore has a neighborhood in which $E_{v,b}$ is
one side of a regular smooth hypersurface.  This proves the claimed
topological and reduced boundary identities, with multiplicity one.
\end{proof}

\subsection{The extra index direction removes the maximum}

Expand the central embedding in the Euclidean parameter
$u=(v,b)\in\mathbb R^p$:
\begin{equation}\label{eq:family-first-variation}
 Y_{v,b}(x)=x+(b+\langle v,x\rangle)N(x)+O(|u|^2).
\end{equation}
The error is smooth and uniform with its source derivatives.
The functions in \eqref{eq:family-normal-variations} are independent:
integrating a relation $t+\langle d,x\rangle=0$ gives $t=0$, and
fullness then gives $d=0$.
To evaluate their second variation, use $L=\Delta+m+P$, $P=|A|^2$,
and write $w=t+f$, $f=\langle d,x\rangle$.
Since $\Delta f=-mf$ and $\int_\Sigma f\,d\mu=0$,
\begin{equation}\label{eq:family-positive-form}
 \int_\Sigma wLw\,d\mu
   =m|\Sigma|t^2+\int_\Sigma P(t+f)^2\,d\mu>0
 \qquad\text{if }(d,t)\ne0.
\end{equation}
For $t\ne0$ the first term suffices.  For $t=0$, $d\ne0$, the
nontrivial analytic function $f$ cannot vanish throughout the nonempty
open set where $P>0$.
The minimal second-variation formula is
$\delta^2\Area(w,w)=-\int_\Sigma wLw\,d\mu$.
Lemma~\ref{lem:family-central-tube} ensures that it computes the mass
of the actual boundary family, proving part~\textup{(iii)} of the
proposition.

Assume now that $\ind(\Sigma)>p$.  Let
\[
 W=\{t+\langle d,x\rangle:t\in\mathbb R,\ d\in\mathbb R^{m+2}\},
 \qquad
 \mathcal B(\phi,\chi)=\int_\Sigma\phi L\chi\,d\mu.
\]
Choose a finite-dimensional subspace spanned by $L$-eigenfunctions
with positive eigenvalues, of dimension greater than $p$.
The $p$ linear conditions
$\mathcal B(\eta,w)=0$ for $w\in W$ admit a nonzero smooth solution
$\eta$ in that subspace.  Then $\mathcal B(\eta,\eta)>0$, and
\eqref{eq:family-positive-form} implies that $\mathcal B$ is positive
definite on $W\oplus\mathbb R\eta$.

Extend $\eta N$ to a smooth ambient vector field on $S^{m+1}$, and
let $F_s$ be its flow.  For small $u,s$, set
\[
 E(u,s)=F_s(E_u),\qquad
 \mathcal A(u,s)=\operatorname{Per}(E(u,s)),
 \qquad E_u=E_{v,b}.
\]
Lemma~\ref{lem:family-central-tube} makes $\mathcal A$ smooth in a neighborhood
of $(0,0)$.  Its first derivative there vanishes by minimality, and
its Hessian is negative definite, by the preceding positive-form
calculation.  Taylor's theorem therefore gives constants $c>0$ and
$\delta>0$ such that
\begin{equation}\label{eq:family-morse-estimate}
 \mathcal A(u,s)\le|\Sigma|-c(|u|^2+s^2)
 \qquad\text{when }|u|^2+s^2\le\delta^2.
\end{equation}
Shrink $\delta$ so that this ball lies inside the central smooth
parameter region.

Choose a smooth nonnegative function $\beta(u)$ supported in
$|u|<\delta/2$ and equal to one near $u=0$.
For sufficiently small $\varepsilon>0$, replace $E_u$ on that region
by $E(u,\varepsilon\beta(u))$ and leave it unchanged elsewhere.
The compact graph
\[
 \{(u,\varepsilon\beta(u)):|u|\le\delta/2\}
\]
misses $(0,0)$.  It is contained in the ball on which
\eqref{eq:family-morse-estimate} holds, so its area maximum is
strictly below $|\Sigma|$.
On the complement of $|u|<\delta/2$, the same strict upper bound,
possibly with a different positive gap, follows from
Lemma~\ref{lem:family-strict-loss}.
The boundaries of the modified domains consequently define a map
$\widetilde\Phi$ satisfying \eqref{eq:family-strict-deformation}.

Replacing $\varepsilon$ by $t\varepsilon$, $0\le t\le1$, gives a
flat homotopy, fixed outside the chosen interior neighborhood of $Q$.
It descends through the quotient and preserves the detecting class.
On the modified region, the boundaries belong to a compact smooth
family of embeddings of $\Sigma$, and hence have no concentration
of mass.  For completeness, a finite cover by coordinate graph
patches, with uniform metric and first-derivative bounds over the
compact parameter set, bounds their masses in radius-$\rho$ balls
by a constant times $\rho^m$.  Outside that region the original
no-concentration estimate applies unchanged.
This proves part~\textup{(iv)}.  Part~\textup{(ii)} is
Lemma~\ref{lem:family-strict-loss}, so the proof of
Proposition~\ref{prop:lorentz-family} is complete.
